\SourceLeaf{159}
Indem man also für die Summe $\alpha A^2+\beta B^2+\gamma C^2+\cdots$ annimmt, vernachlässigt
man alle $(a-A)^2+(a'-A)^2+\cdots$, $(b-B)^2+(b'-B)^2+\cdots$, $\cdots$, die alle positiv sind.

Der Beweis ist leicht: nemlich
\[
\begin{aligned}
aa &= [A+(a-A)]^2 = AA+2A(a-A)+(a-A)^2\\
a'a' &= [A+(a'-A)]^2 = AA+2A(a'-A)+(a'-A)^2\\
a''a'' &= [A+(a''-A)]^2 = AA+2A(a''-A)+(a''-A)^2
\end{aligned}
\]
\begin{center}
. . . . . . . . . . . . . . . . .
\end{center}
\[
\Sigma aa = \alpha AA+2A\Sigma(a-A)+\Sigma(a-A)^2
\]
oder da $\Sigma(a-A)=\Sigma a-\alpha A=0,$
\[
\Sigma aa = \alpha AA+\Sigma(a-A)^2.
\]
\Needspace{6\baselineskip}
\begin{center}
\textsc{Gauss} an \textsc{Olbers}. Göttingen, 3. Mai 1827.
\end{center}

\begin{center}
. . . . . . . . . . . . . . . . .
\end{center}

Zu einer erfolgreichen Anwendung der Wahrscheinlichkeitsrechnung auf Beobachtungen ist
allemal umfassende \textls{Sach}kenntniss von höchster Wichtigkeit. Wo diese fehlt, ist das
Ausschliessen wegen grösserer Differenz immer misslich, wenn nicht die Anzahl der
vorhandenen Beobachtungen sehr gross ist. Alle einzelnen Bestandtheile des
Beobachtungsfehlers, deren Vermeidung ausser unserer Gewalt liegt, haben gewiss Grenzen,
wenn wir auch nicht im Stande sind, sie scharf anzugeben. Es gibt sehr viele Fälle, wo wir
mit Gewissheit sagen können, dass ein vorgekommener grosser Fehler ausserhalb der Grenzen
der Möglichkeit solcher Fehler liegt, und ein ausserordentliches Versehen begangen sein
muss. Die muss man natürlich ausschliessen. So lange man sich aber die Möglichkeit denken
kann, dass der Fehler durch unglückliche Conspiration der Bestandtheile hervorgegangen ist,
soll man nicht ausschliessen. Zuweilen kann es freilich auch Fälle geben, wo man
zweifelhaft ist, ob man sie zur ersten oder zweiten Classe zählen soll; da halte man es wie
man will,
